%; whizzy paragraph -pdf xpdf -latex ./whizzypdfptex.sh
%; whizzy-paragraph "^\\\\begin{frame}\\|\\\\emtext"
% latex beamer presentation.
% platex, latex-beamer でコンパイルすることを想定。 

%     Tokyo Debian Meeting resources
%     Copyright (C) 2012 Junichi Uekawa

%     This program is free software; you can redistribute it and/or modify
%     it under the terms of the GNU General Public License as published by
%     the Free Software Foundation; either version 2 of the License, or
%     (at your option) any later version.

%     This program is distributed in the hope that it will be useful,
%     but WITHOUT ANY WARRANTY; without even the implied warreanty of
%     MERCHANTABILITY or FITNESS FOR A PARTICULAR PURPOSE.  See the
%     GNU General Public License for more details.
%     You should have received a copy of the GNU General Public License
%     along with this program; if not, write to the Free Software
%     Foundation, Inc., 51 Franklin St, Fifth Floor, Boston, MA  02110-1301 USA

\documentclass[cjk,dvipdfmx,12pt]{beamer}
\usetheme{Tokyo}
\usepackage{monthlypresentation}

%  preview (shell-command (concat "evince " (replace-regexp-in-string "tex$" "pdf"(buffer-file-name)) "&")) 
%  presentation (shell-command (concat "xpdf -fullscreen " (replace-regexp-in-string "tex$" "pdf"(buffer-file-name)) "&"))
%  presentation (shell-command (concat "evince " (replace-regexp-in-string "tex$" "pdf"(buffer-file-name)) "&"))

%http://www.naney.org/diki/dk/hyperref.html
%日本語EUC系環境の時
\AtBeginDvi{\special{pdf:tounicode EUC-UCS2}}
%シフトJIS系環境の時
%\AtBeginDvi{\special{pdf:tounicode 90ms-RKSJ-UCS2}}

\newenvironment{commandlinesmall}%
{\VerbatimEnvironment
  \begin{Sbox}\begin{minipage}{1.0\hsize}\begin{fontsize}{8}{8} \begin{BVerbatim}}%
{\end{BVerbatim}\end{fontsize}\end{minipage}\end{Sbox}
  \setlength{\fboxsep}{8pt}
% start on a new paragraph

\vspace{6pt}% skip before
\fcolorbox{dancerdarkblue}{dancerlightblue}{\TheSbox}

\vspace{6pt}% skip after
}
%end of commandlinesmall

\title{東京エリアDebian勉強会}
\subtitle{第146回 2016年11月度}
\author{杉本 典充}
\date{2016年11月19日}
\logo{\includegraphics[width=8cm]{image200607/openlogo-light.eps}}

\begin{document}

\begin{frame}
\titlepage{}
\end{frame}

\begin{frame}{Agenda}
 \begin{minipage}[t]{0.45\hsize}
  \begin{itemize}
  \item 最近あったDebian関連のイベント報告
	\begin{itemize}
	\item 第144回東京エリアDebian勉強会
    \item OSC 2016 Tokyo/Fall (第145回東京エリアDebian勉強会)
	\end{itemize}
  \item 事前課題発表
  \end{itemize}
 \end{minipage}
 \begin{minipage}[t]{0.45\hsize}
  \begin{itemize}
   \item Debian Trivia Quiz
  \end{itemize}
 \end{minipage}
\end{frame}

\section{イベント報告}
\emtext{イベント報告}

\begin{frame}{第144回東京エリアDebian勉強会}

\begin{itemize}
\item 場所は朝日ネットさんをお借りして開催しました
\item 参加者は5名
\item セミナ内容は、以下内容を発表しました
  \begin{enumerate}
  \item dictossさん「preseedでDebianを自動インストールをしてみよう」
  \end{enumerate}
\item 残りの時間でhack timeを行い、成果発表をしました
\item \url{https://tokyodebianmeeting.titanpad.com/ep/pad/export/144/latest?format=txt}
\end{itemize} 
\end{frame}

\begin{frame}{OSC 2016 Tokyo/Fall(第145回)}

\begin{itemize}
\item 11/5(土),11/6(日)に明星大学様で開催されました
\item セミナ内容は、以下内容を発表しました
  \begin{enumerate}
  \item 岩松 信洋さん「Debian updates」
  \end{enumerate}
\item ブースを出展し、OSC参加者へ宣伝・交流・意見交換を行いました
\end{itemize} 
\end{frame}

\section{事前課題}
\emtext{事前課題}
{\footnotesize
 \begin{prework}{ user }
  \begin{enumerate}
  \item xx
  \end{enumerate}
\end{prework}

}

%\subsection{問題}
%
%%; whizzy-master ../debianmeetingresume201311.tex
% 以上の設定をしているため、このファイルで M-x whizzytex すると、whizzytexが利用できます。
%

\santaku
{q}
{A}
{B}
{C}
{B}
{desc}


\section{セミナー1}
\emtext{dh\_strip\_nondeterminism}

\section{Hack time}
\emtext{Hack time}

\begin{frame}{成果を記入下さい!}

  今回、Hack time時の成果を記録に残してみます。

\url{https://tokyodebianmeeting.titanpad.com/146}

に皆さんアクセス頂き（認証不要です）、各自成果を17:30までに記録するようにお願いします。

\end{frame}
  
\section{今後のイベント}
\emtext{今後のイベント}
\begin{frame}{今後のイベント}
  \begin{itemize}
  \item 12/10(土) Mini Debian Conference Japan 2016
    \begin{itemize}
    \item \url{http://miniconf.debian.or.jp/}
    \item 一般参加者、懇親会、KeySignParty参加者募集
    \end{itemize}    
  \item 12月の東京エリアDebian勉強会はお休みの予定です
  \item 2017/1/21(土) 東京エリアDebian勉強会(予定)
  \end{itemize}
\end{frame}

\end{document}

;;; Local Variables: ***
;;; outline-regexp: "\\([ 	]*\\\\\\(documentstyle\\|documentclass\\|emtext\\|section\\|begin{frame}\\)\\*?[ 	]*[[{]\\|[]+\\)" ***
;;; End: ***
